\newpage
\subsubsection{Infraestructura en la nube}
\label{subsec:infraestructura}

La infraestructura del sistema se desplegará sobre Amazon Web Services
(AWS), adoptando un modelo serverless que elimina la necesidad de
gestionar servidores y permite escalar automáticamente según la demanda.
La elección de AWS como proveedor de nube se fundamenta en un conjunto
de factores técnicos, económicos y estratégicos que lo posicionan como
la opción más adecuada para el desarrollo y evolución del prototipo,
tal como se justifica en detalle en la sección~\ref{subsec:justificacion_aws}.

\textbf{Capa gratuita (AWS Free Tier).} AWS ofrece una capa gratuita
que cubre los servicios principales utilizados en este proyecto. A
diferencia de otros proveedores, AWS distingue entre dos tipos de
ofertas gratuitas \cite{aws_free_tier}:

\begin{itemize}
    \item \textbf{Siempre gratuitos (Always Free):} Sin fecha de
    expiración, disponibles indefinidamente mientras se respeten los
    límites mensuales establecidos. Los servicios del sistema cubiertos
    bajo esta categoría son:
    \begin{itemize}
        \item \textbf{AWS Lambda:} 1 millón de invocaciones y 400,000
        GB-segundos de cómputo al mes.
        \item \textbf{Amazon API Gateway:} 1 millón de llamadas REST
        al mes.
        \item \textbf{Amazon CloudFront:} 1 TB de transferencia de
        datos y 10 millones de solicitudes HTTP/S al mes.
        \item \textbf{Amazon S3:} 5 GB de almacenamiento estándar,
        20,000 solicitudes GET y 2,000 solicitudes PUT al mes.
        \item \textbf{Amazon Cognito:} 50,000 usuarios activos mensuales
        sin costo.
    \end{itemize}
    \item \textbf{Gratuitos por 12 meses:} Disponibles desde la
    creación de la cuenta AWS. En esta categoría se encuentra
    \textbf{Amazon RDS}, que ofrece 750 horas mensuales de instancia
    \texttt{db.t2.micro} con PostgreSQL, suficiente para cubrir la
    fase de desarrollo y pruebas del prototipo.
\end{itemize}

Esta disponibilidad permite desarrollar, desplegar y validar el
prototipo completo sin incurrir en costos de infraestructura durante
la fase de desarrollo, lo cual representa una ventaja significativa
para un proyecto académico con recursos limitados.

\textbf{Modelo de pago por uso.} Fuera de la capa gratuita, AWS opera
bajo un modelo de pago por uso donde únicamente se factura por los
recursos efectivamente consumidos, sin compromisos de gasto mínimo
ni contratos de largo plazo. Esto permite que el costo del sistema
crezca de forma proporcional al uso real. Por ejemplo, AWS Lambda
cobra \$0.20 USD por cada millón de invocaciones adicionales y
\$0.0000166667 USD por GB-segundo de cómputo, lo que resulta en
costos mínimos para un prototipo con tráfico reducido
\cite{aws_lambda_pricing}.

\textbf{Facilidad de uso y ecosistema integrado.} AWS provee una
consola de administración unificada, documentación extensa y una
amplia comunidad de soporte. La integración nativa entre sus servicios
reduce considerablemente la complejidad operativa: API Gateway invoca
Lambda directamente, Lambda accede a RDS y S3 dentro de la misma
región sin costos de transferencia adicionales, y CloudFront sirve
el contenido de S3 de forma transparente. Esta cohesión entre servicios
acelera los tiempos de desarrollo y despliegue del sistema.

\textbf{Escalabilidad automática.} El modelo serverless adoptado
mediante AWS Lambda escala automáticamente en función de la demanda,
creando instancias adicionales de forma transparente sin intervención
manual. Este comportamiento es especialmente relevante considerando
que el sistema podría extenderse a otros bosques o reservas naturales
en el futuro, incrementando el número de dispositivos móviles que
sincronizan datos simultáneamente.

\textbf{Disponibilidad global y baja latencia.} AWS cuenta con más
de 30 regiones geográficas y más de 400 puntos de presencia (PoP)
de CloudFront distribuidos globalmente \cite{aws_cloudfront_pricing}.
Para el caso de estudio del Bosque de Chapultepec, la región
\texttt{us-east-1} (Norte de Virginia) ofrece baja latencia desde
México con alta disponibilidad garantizada mediante zonas de
disponibilidad redundantes.

\textbf{Proyección a un ambiente productivo.} El prototipo está
diseñado desde su arquitectura para escalar hacia un ambiente
productivo sin cambios estructurales. En una fase de producción,
el sistema podría incorporar los siguientes servicios adicionales
dentro del mismo ecosistema AWS:

\begin{itemize}
    \item \textbf{Amazon RDS Multi-AZ:} Alta disponibilidad de la
    base de datos mediante replicación automática en múltiples zonas
    de disponibilidad.
    \item \textbf{AWS WAF:} Firewall de aplicaciones web para
    protección contra ataques como inyección SQL y cross-site scripting.
    \item \textbf{Amazon CloudWatch:} Monitoreo en tiempo real de
    métricas, logs y alertas operativas del sistema.
    \item \textbf{Amazon Route 53:} Gestión de dominios y enrutamiento
    de tráfico con alta disponibilidad.
    \item \textbf{AWS Budgets:} Control y alertas de costos para
    evitar gastos inesperados al superar los límites de la capa gratuita.
\end{itemize}

Esta trayectoria de crecimiento es posible gracias a que AWS ofrece
todos estos servicios dentro del mismo ecosistema, sin necesidad de
migrar a otro proveedor ni rediseñar la arquitectura base del sistema.
El análisis comparativo detallado entre AWS, Microsoft Azure y Google
Cloud Platform, así como la justificación económica y técnica de esta
elección, se presentan en la sección~\ref{subsec:justificacion_aws}
del análisis financiero.

A continuación se describen los servicios que conforman la arquitectura
en la nube del sistema.

\textbf{Amazon API Gateway} actuará como punto de entrada unificado
para todas las peticiones REST del sistema, tanto de la aplicación
móvil durante la sincronización como de la aplicación web. Se encargará
de enrutar cada petición hacia la función Lambda correspondiente,
aplicando políticas de control de acceso y limitación de tasa.

\textbf{AWS Lambda} alojará el backend del sistema desarrollado con
Micronaut. Al ser serverless, cada invocación es independiente y el
escalado es automático, eliminando la necesidad de aprovisionar o
gestionar servidores. El código se desplegará automáticamente mediante
el pipeline de CI/CD configurado en GitHub Actions.

\textbf{Amazon RDS (PostgreSQL)} será el motor de base de datos
relacional central del sistema. Almacenará de forma persistente los
datos sincronizados desde los dispositivos móviles, incluyendo
registros de árboles, capturas, resultados de clasificación,
coordenadas GPS e información de usuarios.

\textbf{Amazon S3} gestionará el almacenamiento de objetos no
estructurados del sistema: las imágenes capturadas por los dispositivos
móviles, los artefactos estáticos compilados de la aplicación Angular
y el modelo en formato \texttt{.tflite} distribuido a los dispositivos.

\textbf{Amazon CloudFront} distribuirá globalmente el frontend Angular
mediante una red de distribución de contenido (CDN), sirviendo los
artefactos estáticos almacenados en S3 con baja latencia desde
ubicaciones geográficamente cercanas al usuario.

\textbf{Amazon Cognito} gestionará la identidad y autenticación de
los usuarios de la plataforma web, proveyendo flujos de inicio de
sesión, recuperación de contraseña y emisión de tokens JWT para el
control de acceso por roles.

\textbf{Amazon SageMaker} se empleará exclusivamente como plataforma
de registro y versionado del modelo de clasificación. Permitirá llevar
control de las versiones del modelo entrenado antes de publicarlo en
Amazon S3 para su distribución al dispositivo móvil, sin participar
en el proceso de inferencia en tiempo real.

\textbf{AWS Certificate Manager (ACM)} proveerá y renovará
automáticamente los certificados SSL/TLS asociados a CloudFront y
API Gateway, garantizando que todas las comunicaciones entre los
clientes y la nube se realicen a través de HTTPS, cumpliendo con el
requerimiento no funcional RNF-05.

\textbf{Amazon CloudWatch} se empleará para el monitoreo del sistema
durante el desarrollo y las pruebas del prototipo. Permitirá consultar
los logs generados por las funciones Lambda, monitorear el uso de
Amazon RDS y recibir alertas ante errores o comportamientos inesperados
del backend. Dentro de la capa gratuita, CloudWatch ofrece 5 GB de
ingesta de logs y 10 métricas personalizadas al mes sin costo adicional
\cite{aws_free_tier}.

\textbf{GitHub Actions} automatizará el pipeline de integración y
entrega continua (CI/CD) del backend. Ante cada cambio en el
repositorio, ejecutará las pruebas automatizadas, compilará el
proyecto Micronaut y desplegará la nueva versión en AWS Lambda de
forma transparente.

Adicionalmente, el sistema contará con un \textbf{servidor local
Raspberry Pi} utilizado como entorno de pruebas durante el desarrollo,
permitiendo validar el comportamiento del backend y la sincronización
de datos en un entorno controlado antes del despliegue en la nube.